\pdfoutput=1
\pdfmapfile{+lm.map}
\pdfmapfile{+symbols.map}
\documentclass[11pt]{article}
\usepackage[T1]{fontenc}
\usepackage{lmodern,amsmath,amssymb,amsthm,hyperref}
\usepackage[margin=1in]{geometry}
\newtheorem{proposition}{Proposition}
\title{The tautological-to-Chow inclusion has zero kernel\\A disproof of Conjecture 00000007811}
\author{Galaxy-0}\date{October 19, 2026}
\begin{document}\maketitle
\begin{abstract}
The nontrivial-kernel clause in Conjecture 7811 contradicts the definition of the Chow tautological ring as a subring of the Chow ring. Its natural map into the Chow ring is inclusion, hence injective in every genus. We give the complete argument, distinguish this map from a presentation map and a cycle-class map, and supply a universally quantified Lean proof of the relevant algebraic fact with no axioms.
\end{abstract}
\section{Statement and meaning of the map}
The conjecture asserts, in particular, that the natural map
\[
R^*(\mathcal M_g)\longrightarrow A^*(\mathcal M_g)
\]
has nontrivial kernel for $g\ge8$, with its first nontrivial relation in degree $g-4$.
Here $A^*$ is the Chow ring with rational coefficients. In this setting the tautological ring $R^*(\mathcal M_g)$ is, by definition, the subring of $A^*(\mathcal M_g)$ generated by the tautological classes. Whether one lists only the $\kappa$ classes or also the $\lambda$ classes makes no difference to the following argument: a subring generated by any classes is still a subring of the ambient ring.

This convention is explicit in Vakil, \emph{The moduli space of curves and Gromov-Witten theory}, Section 3.1, which defines $R^*(\mathcal M_g)\subset A^*(\mathcal M_g)$; see \url{https://math.stanford.edu/~vakil/files/cimefeb19.pdf}. It is also stated in Section 2.2 of Pixton, \emph{The tautological ring of the moduli space of curves}, \url{https://websites.umich.edu/~pixton/papers/tautological.pdf}. The natural map to the same Chow ring is the inclusion, not a map from formal symbols before imposing relations.

\section{Disproof}
\begin{proposition}
For every ring $A$ and every subring $R\subseteq A$, the inclusion $\iota:R\to A$ has kernel $\{0\}$.
\end{proposition}
\begin{proof}
An element $r$ of $R$ is an element of $A$ satisfying the subring membership condition. The inclusion sends this element to itself. Consequently $\iota(r)=0_A$ implies $r=0_R$; conversely the zero element is sent to zero. This proves the assertion.
\end{proof}
Apply the proposition with $A=A^*(\mathcal M_8)$ and $R=R^*(\mathcal M_8)$. The genus satisfies $8\ge8$, but the kernel is zero, contrary to the asserted nontriviality. Indeed the kernel is zero for every genus and in every degree, including degree $8-4=4$. Therefore the conjunction in Conjecture 7811 is false. No assertion about the conjecture's separate alternating-relations/cohomology clause is needed.

\section{Relations do not give counterexamples to injectivity}
If $S$ is a polynomial or strata algebra on formal tautological generators, the evaluation map $S\twoheadrightarrow R^*$ can have a nonzero kernel: its elements are presentation relations. After passing to $R^*$ these relations are already zero. They do not become nonzero elements of the kernel of $R^*\hookrightarrow A^*$.
Likewise, the cycle-class map from Chow classes to cohomology has a different target and is not the map printed in the conjecture. If one instead uses the cohomological tautological ring, a natural reverse map into the Chow ring is not provided by that convention. Our disproof concerns the stated tautological-to-Chow map under its standard Chow-theoretic meaning, not a conjecture about either of these different maps. It does not refute Faber's Gorenstein conjecture.

\section{Formalization and its scope}
The self-contained Lean 4.31.0 project uses only the standard library. It implements an ambient carrier with zero, one, addition, multiplication and negation, then defines \texttt{Generated} inductively by closure under these operations and the specified generators. The resulting \texttt{GeneratedSubring} consists of ambient elements with a proof of membership. Including all rational scalar constants among the generators covers generation as a rational subalgebra.

Ring laws are not assumed or needed: the theorem holds for every carrier with these operations, and therefore for all rings and all Chow rings. Injectivity is proved from subtype equality, not inserted as a hypothesis. The following results are checked:
\begin{itemize}
\item \texttt{inclusion\_injective}: equality of ambient images implies equality in the generated subring.
\item \texttt{kernel\_iff\_zero} and \texttt{no\_nontrivial\_kernel}: the kernel consists exactly of zero.
\item \texttt{genus\_kernel\_claim\_false}: for arbitrary ambient algebras and generating classes indexed by the genus, the universal nontrivial-kernel claim for genus at least eight is false, by specialization at eight.
\item \texttt{no\_kernel\_in\_any\_degree}: adding any condition specifying a degree cannot produce a nonzero kernel element.
\end{itemize}
The code does not construct the moduli stack, compute its Chow ring, or prove results about the alternating cohomology clause. Such constructions are unnecessary for the universal injectivity argument. The semantic link to the geometric notation is exactly the standard definition cited above; no claim is made that those geometric constructions have been formalized.

\section{Reproduction}
From \texttt{lean/}, with Lean 4.31.0 on the path:
\begin{verbatim}
lake build
lake env lean -DwarningAsError=true Main.lean
lake env lean Audit.lean
lake env leanchecker Main
\end{verbatim}
The five audited theorems have no axiom dependencies. There are no proof holes, custom axioms, native decision procedures or kernel bypasses. The report can be built with \texttt{pdflatex -no-shell-escape solution.tex}. No auxiliary numerical computation is needed.
\end{document}
